\section{Reduce an arbitrary input a bit at a time}
Modular addition expects residues. The public multiplication and power interfaces
accept arbitrary unsigned 128-bit inputs, so they need a reduction operation first.
One approach is binary long division: consume the input from its highest bit down,
retaining only the remainder of the prefix already seen.

Start \(R=0\). For each input bit \(b\), replace \(R\) with \((2R+b)\bmod N\).
The new binary prefix has exactly this value before taking its remainder. Implement
its double with \texttt{rsa\_addmod128}(\(R,R,N,R\)). If the input bit is one, add
a private object representing one with a second modular addition. Because the addition
interface permits exact aliasing, both calls can update the private remainder in place.

For the four bits of decimal 13, reducing modulo 5 gives:
\tagpdfsetup{table/header-rows={1}}
\begin{center}
  \begin{tabular}{rrrr}
    Next bit & Old remainder & \(2R+b\) & New remainder \\
    \hline
    1        & 0             & 1        & 1             \\
    1        & 1             & 3        & 3             \\
    0        & 3             & 6        & 1             \\
    1        & 1             & 3        & 3             \\
  \end{tabular}
\end{center}
Leading zero bits leave the initial remainder zero. The 128-bit implementation can
therefore use a fixed scan from bit 127 down to zero. An input already below \(N\)
can instead be copied directly after an unsigned multiword comparison.

To find bit \(k\), select word \(\lfloor k/32\rfloor\), shift it logically by
\(k\bmod32\), and mask with one. In this fragment \texttt{s0} is the original
input pointer and \texttt{s3} is \(k\). The pointer addition is doubleword-sized.

\Needspace{13\baselineskip}
\MARSCodeLabel{Extract one bit without confusing a bit index with a byte offset.}
\begin{minted}[samepage=true]{mars-idea}
dsrl $t0, $s3, 5
    dsll $t0, $t0, 2
    daddu $t0, $s0, $t0
    lwu $t1, 0($t0)
    andi $t2, $s3, 31
    dsrlv $t1, $t1, $t2
    andi $t1, $t1, 1
\end{minted}

Handle modulus one separately by writing zero. A private object containing one would
not be a valid residue for the modular-addition precondition when \(N=1\). For \(N>1\),
both zero and one are valid residues. Handle modulus zero first, before either input
inspection or output stores. These special cases let the loop use its invariant
without adding exceptions to each nested call.
